% =====================================================================
% Note on: CDS Market Formulas and Models
% Damiano Brigo and Massimo Morini, 2005
% =====================================================================
\documentclass[11pt,a4paper]{article}

\newcommand{\papertag}{brigo\_morini\_2005\_cdsmm}

% =====================================================================
% Shared preamble for paper notes
% Include with:  % =====================================================================
% Shared preamble for paper notes
% Include with:  % =====================================================================
% Shared preamble for paper notes
% Include with:  \input{../_common/preamble}
% =====================================================================

% --- Core packages ---
\usepackage[utf8]{inputenc}
\usepackage[T1]{fontenc}
\usepackage{lmodern}
\usepackage[margin=2.5cm]{geometry}
\usepackage{amsmath,amssymb,amsthm,mathtools}
\usepackage{bm}
\usepackage{booktabs}
\usepackage{array}
\usepackage{enumitem}
\usepackage{hyperref}
\usepackage{xcolor}
\usepackage{listings}
\usepackage{fancyhdr}
\usepackage{titlesec}

% --- Hyperref styling ---
\hypersetup{
  colorlinks=true,
  linkcolor=blue!50!black,
  urlcolor=blue!50!black,
  citecolor=blue!50!black
}

% --- Theorem-like environments ---
\theoremstyle{definition}
\newtheorem{defn}{Definition}
\newtheorem{remark}{Remark}
\newtheorem{example}{Example}

\theoremstyle{plain}
\newtheorem{prop}{Proposition}
\newtheorem{lemma}{Lemma}

% --- Coloured boxes for the structured sections ---
% Validation block, commentary, TODOs use coloured frames so the eye
% can find them when flipping through the note.
\usepackage[most]{tcolorbox}

\newtcolorbox{commentary}{
  colback=yellow!5,
  colframe=yellow!50!black,
  title=Commentary,
  fonttitle=\bfseries,
  breakable
}

\newtcolorbox{validation}{
  colback=green!3,
  colframe=green!50!black,
  title=Validation,
  fonttitle=\bfseries,
  breakable
}

\newtcolorbox{todobox}{
  colback=red!3,
  colframe=red!50!black,
  title=TODO / Open questions,
  fonttitle=\bfseries,
  breakable
}

% --- Code listing style for Python snippets in validation blocks ---
\lstdefinestyle{py}{
  language=Python,
  basicstyle=\ttfamily\footnotesize,
  keywordstyle=\color{blue!60!black},
  commentstyle=\color{gray},
  stringstyle=\color{red!60!black},
  showstringspaces=false,
  breaklines=true,
  frame=single,
  framesep=4pt,
  rulecolor=\color{gray!30}
}

% --- Page header: shows paper tag so a printed note is identifiable ---
\pagestyle{fancy}
\fancyhf{}
\fancyhead[L]{\small\textit{\papertag}}
\fancyhead[R]{\small\thepage}
\renewcommand{\headrulewidth}{0.4pt}

% --- Section spacing: tighter than article default ---
\titlespacing*{\section}{0pt}{1.5ex plus 0.5ex minus 0.2ex}{1ex plus 0.2ex}
\titlespacing*{\subsection}{0pt}{1.2ex plus 0.3ex minus 0.2ex}{0.6ex plus 0.2ex}

% =====================================================================

% --- Core packages ---
\usepackage[utf8]{inputenc}
\usepackage[T1]{fontenc}
\usepackage{lmodern}
\usepackage[margin=2.5cm]{geometry}
\usepackage{amsmath,amssymb,amsthm,mathtools}
\usepackage{bm}
\usepackage{booktabs}
\usepackage{array}
\usepackage{enumitem}
\usepackage{hyperref}
\usepackage{xcolor}
\usepackage{listings}
\usepackage{fancyhdr}
\usepackage{titlesec}

% --- Hyperref styling ---
\hypersetup{
  colorlinks=true,
  linkcolor=blue!50!black,
  urlcolor=blue!50!black,
  citecolor=blue!50!black
}

% --- Theorem-like environments ---
\theoremstyle{definition}
\newtheorem{defn}{Definition}
\newtheorem{remark}{Remark}
\newtheorem{example}{Example}

\theoremstyle{plain}
\newtheorem{prop}{Proposition}
\newtheorem{lemma}{Lemma}

% --- Coloured boxes for the structured sections ---
% Validation block, commentary, TODOs use coloured frames so the eye
% can find them when flipping through the note.
\usepackage[most]{tcolorbox}

\newtcolorbox{commentary}{
  colback=yellow!5,
  colframe=yellow!50!black,
  title=Commentary,
  fonttitle=\bfseries,
  breakable
}

\newtcolorbox{validation}{
  colback=green!3,
  colframe=green!50!black,
  title=Validation,
  fonttitle=\bfseries,
  breakable
}

\newtcolorbox{todobox}{
  colback=red!3,
  colframe=red!50!black,
  title=TODO / Open questions,
  fonttitle=\bfseries,
  breakable
}

% --- Code listing style for Python snippets in validation blocks ---
\lstdefinestyle{py}{
  language=Python,
  basicstyle=\ttfamily\footnotesize,
  keywordstyle=\color{blue!60!black},
  commentstyle=\color{gray},
  stringstyle=\color{red!60!black},
  showstringspaces=false,
  breaklines=true,
  frame=single,
  framesep=4pt,
  rulecolor=\color{gray!30}
}

% --- Page header: shows paper tag so a printed note is identifiable ---
\pagestyle{fancy}
\fancyhf{}
\fancyhead[L]{\small\textit{\papertag}}
\fancyhead[R]{\small\thepage}
\renewcommand{\headrulewidth}{0.4pt}

% --- Section spacing: tighter than article default ---
\titlespacing*{\section}{0pt}{1.5ex plus 0.5ex minus 0.2ex}{1ex plus 0.2ex}
\titlespacing*{\subsection}{0pt}{1.2ex plus 0.3ex minus 0.2ex}{0.6ex plus 0.2ex}

% =====================================================================

% --- Core packages ---
\usepackage[utf8]{inputenc}
\usepackage[T1]{fontenc}
\usepackage{lmodern}
\usepackage[margin=2.5cm]{geometry}
\usepackage{amsmath,amssymb,amsthm,mathtools}
\usepackage{bm}
\usepackage{booktabs}
\usepackage{array}
\usepackage{enumitem}
\usepackage{hyperref}
\usepackage{xcolor}
\usepackage{listings}
\usepackage{fancyhdr}
\usepackage{titlesec}

% --- Hyperref styling ---
\hypersetup{
  colorlinks=true,
  linkcolor=blue!50!black,
  urlcolor=blue!50!black,
  citecolor=blue!50!black
}

% --- Theorem-like environments ---
\theoremstyle{definition}
\newtheorem{defn}{Definition}
\newtheorem{remark}{Remark}
\newtheorem{example}{Example}

\theoremstyle{plain}
\newtheorem{prop}{Proposition}
\newtheorem{lemma}{Lemma}

% --- Coloured boxes for the structured sections ---
% Validation block, commentary, TODOs use coloured frames so the eye
% can find them when flipping through the note.
\usepackage[most]{tcolorbox}

\newtcolorbox{commentary}{
  colback=yellow!5,
  colframe=yellow!50!black,
  title=Commentary,
  fonttitle=\bfseries,
  breakable
}

\newtcolorbox{validation}{
  colback=green!3,
  colframe=green!50!black,
  title=Validation,
  fonttitle=\bfseries,
  breakable
}

\newtcolorbox{todobox}{
  colback=red!3,
  colframe=red!50!black,
  title=TODO / Open questions,
  fonttitle=\bfseries,
  breakable
}

% --- Code listing style for Python snippets in validation blocks ---
\lstdefinestyle{py}{
  language=Python,
  basicstyle=\ttfamily\footnotesize,
  keywordstyle=\color{blue!60!black},
  commentstyle=\color{gray},
  stringstyle=\color{red!60!black},
  showstringspaces=false,
  breaklines=true,
  frame=single,
  framesep=4pt,
  rulecolor=\color{gray!30}
}

% --- Page header: shows paper tag so a printed note is identifiable ---
\pagestyle{fancy}
\fancyhf{}
\fancyhead[L]{\small\textit{\papertag}}
\fancyhead[R]{\small\thepage}
\renewcommand{\headrulewidth}{0.4pt}

% --- Section spacing: tighter than article default ---
\titlespacing*{\section}{0pt}{1.5ex plus 0.5ex minus 0.2ex}{1ex plus 0.2ex}
\titlespacing*{\subsection}{0pt}{1.2ex plus 0.3ex minus 0.2ex}{0.6ex plus 0.2ex}

% =====================================================================
% House notation: shared symbols used across all paper notes.
% Each note imports this and may shadow individual macros if a paper
% uses a clashing convention — but the *default* meaning is fixed here.
%
% Include with:  % =====================================================================
% House notation: shared symbols used across all paper notes.
% Each note imports this and may shadow individual macros if a paper
% uses a clashing convention — but the *default* meaning is fixed here.
%
% Include with:  % =====================================================================
% House notation: shared symbols used across all paper notes.
% Each note imports this and may shadow individual macros if a paper
% uses a clashing convention — but the *default* meaning is fixed here.
%
% Include with:  \input{../_common/notation}
% =====================================================================

% --- Measures and expectations ---
\newcommand{\Q}{\mathbb{Q}}              % risk-neutral measure
\newcommand{\Qt}{\mathbb{Q}^{T}}         % T-forward measure
\newcommand{\E}{\mathbb{E}}              % expectation
\newcommand{\Et}[1]{\E^{#1}}             % expectation under measure #1
\newcommand{\Ftime}[1]{\mathcal{F}_{#1}} % filtration at #1

% --- Discount and forward objects ---
\newcommand{\Disc}[2]{D_{#1,#2}}         % discount factor D_{t,T}
\newcommand{\Bond}[2]{P_{#1}(#2)}        % bond price P_t(y)
\newcommand{\Ann}[1]{A_{#1}}             % annuity / level
\newcommand{\Numer}[1]{N_{#1}}           % numeraire
\newcommand{\Fwd}[2]{F_{#1,#2}}          % forward price F_{t,T}
\newcommand{\Fwdpx}[2]{\mathrm{Fwd}_{#1}(#2)} % verbose forward-price F_t(T)

% --- Rates ---
\newcommand{\IRR}{\mathrm{IRR}}          % internal rate of return
\newcommand{\Yld}{y}                     % bond yield variable
\newcommand{\Strike}{K}                  % strike (price-space)
\newcommand{\Lock}{L}                    % locked yield / rate
\newcommand{\Repo}{r^{\mathrm{repo}}}    % repo rate
\newcommand{\Fund}{r^{\mathrm{fun}}}     % unsecured funding rate
\newcommand{\OIS}{r^{\mathrm{ois}}}      % OIS rate
\newcommand{\Coll}{r^{\mathrm{c}}}       % collateral rate
\newcommand{\Rcmt}{R^{\mathrm{cmt}}}     % constant-maturity-treasury rate
\newcommand{\Rswp}{R^{\mathrm{swp}}}     % swap rate (used for risky variant)
\newcommand{\Rec}{\mathrm{Rec}}          % recovery rate
\newcommand{\NPV}{\mathrm{NPV}}          % present-value functional
\newcommand{\PVone}{\mathrm{PV01}}       % risky annuity
\newcommand{\Loss}{\mathrm{Loss}}        % default-leg integral
\newcommand{\Sasw}{S^{\mathrm{ASW}}}     % asset-swap spread
\newcommand{\Scds}{S^{\mathrm{CDS}}}     % CDS spread
\newcommand{\Basis}{\mathrm{Basis}}      % CDS-bond basis

% --- Sensitivities ---
\newcommand{\Diff}[2]{D_{#1}^{#2}}       % Euler's notation: D_x^k[f]
\newcommand{\Riskf}{\mathrm{RiskFactor}} % bond risk factor (= -DV01*1e4)
\newcommand{\DVone}{\mathrm{DV01}}

% --- Indicator / misc ---
\newcommand{\Ind}{\mathbf{1}}
\newcommand{\given}{\,|\,}
\newcommand{\dd}{\,\mathrm{d}}

% --- Greeks-style ---
\newcommand{\Gam}{\Gamma}
\newcommand{\Vega}{\mathcal{V}}

% =====================================================================

% --- Measures and expectations ---
\newcommand{\Q}{\mathbb{Q}}              % risk-neutral measure
\newcommand{\Qt}{\mathbb{Q}^{T}}         % T-forward measure
\newcommand{\E}{\mathbb{E}}              % expectation
\newcommand{\Et}[1]{\E^{#1}}             % expectation under measure #1
\newcommand{\Ftime}[1]{\mathcal{F}_{#1}} % filtration at #1

% --- Discount and forward objects ---
\newcommand{\Disc}[2]{D_{#1,#2}}         % discount factor D_{t,T}
\newcommand{\Bond}[2]{P_{#1}(#2)}        % bond price P_t(y)
\newcommand{\Ann}[1]{A_{#1}}             % annuity / level
\newcommand{\Numer}[1]{N_{#1}}           % numeraire
\newcommand{\Fwd}[2]{F_{#1,#2}}          % forward price F_{t,T}
\newcommand{\Fwdpx}[2]{\mathrm{Fwd}_{#1}(#2)} % verbose forward-price F_t(T)

% --- Rates ---
\newcommand{\IRR}{\mathrm{IRR}}          % internal rate of return
\newcommand{\Yld}{y}                     % bond yield variable
\newcommand{\Strike}{K}                  % strike (price-space)
\newcommand{\Lock}{L}                    % locked yield / rate
\newcommand{\Repo}{r^{\mathrm{repo}}}    % repo rate
\newcommand{\Fund}{r^{\mathrm{fun}}}     % unsecured funding rate
\newcommand{\OIS}{r^{\mathrm{ois}}}      % OIS rate
\newcommand{\Coll}{r^{\mathrm{c}}}       % collateral rate
\newcommand{\Rcmt}{R^{\mathrm{cmt}}}     % constant-maturity-treasury rate
\newcommand{\Rswp}{R^{\mathrm{swp}}}     % swap rate (used for risky variant)
\newcommand{\Rec}{\mathrm{Rec}}          % recovery rate
\newcommand{\NPV}{\mathrm{NPV}}          % present-value functional
\newcommand{\PVone}{\mathrm{PV01}}       % risky annuity
\newcommand{\Loss}{\mathrm{Loss}}        % default-leg integral
\newcommand{\Sasw}{S^{\mathrm{ASW}}}     % asset-swap spread
\newcommand{\Scds}{S^{\mathrm{CDS}}}     % CDS spread
\newcommand{\Basis}{\mathrm{Basis}}      % CDS-bond basis

% --- Sensitivities ---
\newcommand{\Diff}[2]{D_{#1}^{#2}}       % Euler's notation: D_x^k[f]
\newcommand{\Riskf}{\mathrm{RiskFactor}} % bond risk factor (= -DV01*1e4)
\newcommand{\DVone}{\mathrm{DV01}}

% --- Indicator / misc ---
\newcommand{\Ind}{\mathbf{1}}
\newcommand{\given}{\,|\,}
\newcommand{\dd}{\,\mathrm{d}}

% --- Greeks-style ---
\newcommand{\Gam}{\Gamma}
\newcommand{\Vega}{\mathcal{V}}

% =====================================================================

% --- Measures and expectations ---
\newcommand{\Q}{\mathbb{Q}}              % risk-neutral measure
\newcommand{\Qt}{\mathbb{Q}^{T}}         % T-forward measure
\newcommand{\E}{\mathbb{E}}              % expectation
\newcommand{\Et}[1]{\E^{#1}}             % expectation under measure #1
\newcommand{\Ftime}[1]{\mathcal{F}_{#1}} % filtration at #1

% --- Discount and forward objects ---
\newcommand{\Disc}[2]{D_{#1,#2}}         % discount factor D_{t,T}
\newcommand{\Bond}[2]{P_{#1}(#2)}        % bond price P_t(y)
\newcommand{\Ann}[1]{A_{#1}}             % annuity / level
\newcommand{\Numer}[1]{N_{#1}}           % numeraire
\newcommand{\Fwd}[2]{F_{#1,#2}}          % forward price F_{t,T}
\newcommand{\Fwdpx}[2]{\mathrm{Fwd}_{#1}(#2)} % verbose forward-price F_t(T)

% --- Rates ---
\newcommand{\IRR}{\mathrm{IRR}}          % internal rate of return
\newcommand{\Yld}{y}                     % bond yield variable
\newcommand{\Strike}{K}                  % strike (price-space)
\newcommand{\Lock}{L}                    % locked yield / rate
\newcommand{\Repo}{r^{\mathrm{repo}}}    % repo rate
\newcommand{\Fund}{r^{\mathrm{fun}}}     % unsecured funding rate
\newcommand{\OIS}{r^{\mathrm{ois}}}      % OIS rate
\newcommand{\Coll}{r^{\mathrm{c}}}       % collateral rate
\newcommand{\Rcmt}{R^{\mathrm{cmt}}}     % constant-maturity-treasury rate
\newcommand{\Rswp}{R^{\mathrm{swp}}}     % swap rate (used for risky variant)
\newcommand{\Rec}{\mathrm{Rec}}          % recovery rate
\newcommand{\NPV}{\mathrm{NPV}}          % present-value functional
\newcommand{\PVone}{\mathrm{PV01}}       % risky annuity
\newcommand{\Loss}{\mathrm{Loss}}        % default-leg integral
\newcommand{\Sasw}{S^{\mathrm{ASW}}}     % asset-swap spread
\newcommand{\Scds}{S^{\mathrm{CDS}}}     % CDS spread
\newcommand{\Basis}{\mathrm{Basis}}      % CDS-bond basis

% --- Sensitivities ---
\newcommand{\Diff}[2]{D_{#1}^{#2}}       % Euler's notation: D_x^k[f]
\newcommand{\Riskf}{\mathrm{RiskFactor}} % bond risk factor (= -DV01*1e4)
\newcommand{\DVone}{\mathrm{DV01}}

% --- Indicator / misc ---
\newcommand{\Ind}{\mathbf{1}}
\newcommand{\given}{\,|\,}
\newcommand{\dd}{\,\mathrm{d}}

% --- Greeks-style ---
\newcommand{\Gam}{\Gamma}
\newcommand{\Vega}{\mathcal{V}}


\title{Note on: \emph{CDS Market Formulas and Models}\\
       \large Damiano Brigo \& Massimo Morini, 2005}
\author{Reading notes}
\date{\today}

\begin{document}
\maketitle

% =====================================================================
\section*{Metadata}
\begin{tabular}{@{}p{3.5cm}p{11cm}@{}}
\toprule
Title           & CDS Market Formulas and Models \\
Authors         & Damiano Brigo (Banca IMI), Massimo Morini (Università di Milano Bicocca) \\
Date            & September 2005 \\
Source          & Working paper, no SSRN ID printed on document \\
Local file      & \texttt{papers/brigo\_morini\_2005\_cdsmm/cdsmktfor.pdf} \\
Tags            & \texttt{cds}, \texttt{cds-option}, \texttt{market-model}, \texttt{cmcds},
                  \texttt{subfiltration}, \texttt{change-of-measure}, \texttt{black-formula} \\
Date read       & \today \\
\bottomrule
\end{tabular}

% =====================================================================
\section*{One-paragraph summary}
The paper builds a Swap Market Model analogue for credit derivatives:
a CDS Market Model in which the par CDS spread $R_{a,b}(t)$ is the
natural underlying, with a Black-Scholes-type closed-form pricing
formula for European CDS options under a lognormal-spread assumption.
The construction uses the Jeanblanc-Rutkowski subfiltration setup ---
$\mathcal{F}_t = \mathcal{F}^\tau_t \vee \mathcal{H}_t$, separating
default information from default-free information --- to keep all
measures equivalent to the risk-neutral measure (avoiding the
defaultable-bond numeraire pathology in Schonbucher 1999). The
defaultable PV01 $C_{a,b}(t) = \sum_i \alpha_i \bar P(t, T_i)$ serves
as numeraire, defining a measure $\bar Q^{a,b}$ under which the par
spread is a martingale, yielding the Black formula
\eqref{eq:cdsmm-black}. Under the conditional-independence assumption
between the subfiltration $\mathcal{H}_t$ and the risk-neutral
measure, the framework also allows extension to a CDS term-structure
model by parameterising one-period rates $R_j(t)$, with the
practical complication that defaultable-bond numeraire ratios cannot
be expressed purely in terms of $R_j$ unless one further assumes
independence of default-free interest rates and default probabilities.
Under this independence, a full LMM-style CDS market model is
defined and applied to constant-maturity CDS (CMCDS) via the
drift-freezing approximation; a closed-form CMCDS price is obtained.
Empirical examples on European corporate CDS options (Deutsche
Telecom, Daimler Chrysler, France Telecom) and a CMCDS on Fiat are
included.

% =====================================================================
\section{Setup and notation}

\begin{tabular}{@{}p{3cm}p{3.5cm}p{8.5cm}@{}}
\toprule
Paper                       & House                          & Meaning \\
\midrule
$\tau$                      & $\tau$                         & Default time of reference entity \\
$T_a, \ldots, T_b$          & $T_i$                          & Payment dates of the CDS premium leg \\
$\alpha_i$                  & $\alpha_i$                     & Year fraction $T_{i-1} \to T_i$ \\
$\mathrm{LGD}$              & $1 - \Rec$                     & Loss given default \\
$R$                         & $R$                            & Running CDS premium / spread paid by buyer \\
$R_{a,b}(t)$                & $\Scds(t)$                     & Par (fair) CDS spread, maturity $T_a \to T_b$ \\
$D(t, T)$                   & $\Disc{t}{T}$                  & Stochastic bank-account discount $B_t / B_T$ \\
$P(t, T)$                   & $\Disc{t}{T}$                  & Default-free zero-coupon bond \\
$\bar P(t, T)$              & ---                            & Defaultable zero (pre-default) $= \E^Q[D(t,T)\,\Ind_{\tau > T} | \mathcal{H}_t]/Q(\tau > t|\mathcal{H}_t)$ \\
$Q(\tau > t | \mathcal{H}_t)$ & $Q(t, T)$                    & Conditional survival probability \\
$\mathcal{F}_t$             & ---                            & Full filtration (default time visible) \\
$\mathcal{H}_t$             & ---                            & Default-free subfiltration \\
$\mathcal{F}^\tau_t$        & ---                            & Subfiltration generated by $\tau$ \\
$Q$                         & $\Q$                           & Risk-neutral measure \\
$\bar Q^{a,b}$              & ---                            & Measure with numeraire $C_{a,b}$ (defaultable PV01) \\
$C_{a,b}(t)$                & $\PVone^{\text{Risky}}_{a,b}(t)$ & Defaultable PV01 $\sum_i \alpha_i Q(\tau{>}t|\mathcal{H}_t)\bar P(t, T_i)$ \\
$Q^j$                       & ---                            & $T_j$-forward measure of default-free bond $P(t, T_j)$ \\
$Q^j(\tau > T | \mathcal{H}_t)$ & ---                        & Forward conditional survival \\
$F_k(t)$                    & ---                            & Default-free forward Libor (LMM variable) \\
$R_j(t) = R_{j-1,j}(t)$     & ---                            & One-period forward CDS spread \\
$\bar\sigma_{a,b}$          & ---                            & Implied volatility of $R_{a,b}$ under $\bar Q^{a,b}$ \\
$\bar{\text{Black}}(F, K, v)$ & ---                          & Standard Black formula \\
\bottomrule
\end{tabular}

\paragraph{Subfiltration setup.} The full $\sigma$-algebra is
$\mathcal{F}_t = \mathcal{F}^\tau_t \vee \mathcal{H}_t$: full
information equals knowledge of the default time plus knowledge of
all default-free observables (rates, market quotes). The point is
that the conditional survival probability $Q(\tau > t | \mathcal{H}_t)$
is strictly positive in every state of the world, while
$Q(\tau > t | \mathcal{F}_t) = \Ind_{\tau > t}$ jumps to zero at
default. Working under $\mathcal{H}_t$ keeps the numeraire $C_{a,b}$
strictly positive, hence all change-of-measure machinery works.

% =====================================================================
\section{Key equations}

\paragraph{CDS payoff convention} (paper's eq.~1, slightly more
tractable than the textbook payment-at-default):
\begin{equation}\label{eq:cdsmm-payoff}
  \text{CDS}\Pi_t(R) \;=\; \mathrm{LGD} \sum_{i=a+1}^{b} D(t, T_i)\, \Ind_{T_{i-1} < \tau \le T_i}
                       \;-\; \sum_{i=a+1}^{b} D(t, T_i)\, \alpha_i\, \Ind_{\tau > T_i}\, R.
\end{equation}
\textit{Says:} the protection leg pays $\mathrm{LGD}$ at the first
$T_i$ following default (not at $\tau$ itself --- the
``postponed-default'' simplification); the premium leg pays $R$
times $\alpha_i$ at each $T_i$ if alive, with \emph{no} accrued
premium on the default-period stub.

\paragraph{Par CDS spread} (paper's central definition, with
subfiltration):
\begin{equation}\label{eq:cdsmm-par}
  R_{a,b}(t) \;=\; \mathrm{LGD} \cdot
                \frac{\sum_{i=a+1}^{b} \E^Q\!\bigl[D(t, T_i)\,\Ind_{T_{i-1} < \tau \le T_i}\,|\,\mathcal{H}_t\bigr]}
                     {\sum_{i=a+1}^{b} \alpha_i\, Q(\tau > t | \mathcal{H}_t)\, \bar P(t, T_i)},
\end{equation}
with $\bar P(t, T) = \E^Q[D(t,T)\Ind_{\tau > T}|\mathcal{H}_t]/Q(\tau > t|\mathcal{H}_t)$
the pre-default defaultable bond price.
\textit{Says:} the par spread is a ratio of expected protection
payments to the defaultable annuity. Using $\mathcal{H}_t$
conditioning rather than $\mathcal{F}_t$ avoids zero denominators at
default.

\paragraph{Defaultable PV01 numeraire:}
\begin{equation}\label{eq:cdsmm-pv01}
  C_{a,b}(t) \;:=\; \sum_{i=a+1}^{b} \alpha_i\, Q(\tau > t | \mathcal{H}_t)\, \bar P(t, T_i)
            \;=\; \sum_{i=a+1}^{b} \alpha_i\, \E^Q\!\bigl[D(t, T_i)\,\Ind_{\tau > T_i}\,|\,\mathcal{H}_t\bigr].
\end{equation}
\textit{Says:} the defaultable PV01 (``DV01 of one bp of premium'')
is strictly positive (Proposition~7 of the paper); hence its
accreted version $C_{a,b} = C_{a,b}(T_a)\, B/B_{T_a}$ is a valid
numeraire, defining the measure $\bar Q^{a,b}$ equivalent to $Q$.

\paragraph{Standard Market Formula for CDS options} (paper's eq.~8):
\begin{equation}\label{eq:cdsmm-black}
\boxed{\;
  \text{CDSOption}_t \;=\; \Ind_{\tau > t} \sum_{i=a+1}^{b}
          \alpha_i\, \bar P(t, T_i)\,
          \text{Black}\!\bigl(R_{a,b}(t),\, K,\, \bar\sigma_{a,b}\sqrt{T_a - t}\bigr),\;
}
\end{equation}
with $\text{Black}(F, K, v) = F\, N(d_1) - K\, N(d_2)$,
$d_{1,2} = [\ln(F/K) \pm \tfrac12 v^2]/v$, under the lognormal
assumption $\dd R_{a,b}(t) = \bar\sigma_{a,b}\, R_{a,b}(t)\, \dd V^{a,b}$
where $V^{a,b}$ is Brownian motion under $\bar Q^{a,b}$.
\textit{Says:} the CDS option prices like a Black caplet/swaption,
with the par CDS spread as the underlying and the defaultable PV01
as the discounted notional. The novelty is that, thanks to the
subfiltration setup, $\bar Q^{a,b}$ is equivalent to $Q$ (not
absolutely continuous only), so implied vols are well-defined.

\paragraph{CDS spread as forward survival-probability ratio}
(paper's eq.~13, central one-period structural identity, under
$\mathcal{H}$-conditional independence):
\begin{equation}\label{eq:cdsmm-onerate}
  R_j(t) \;=\; \frac{\mathrm{LGD}}{\alpha_j}
              \left( \frac{Q^j(\tau > T_{j-1} | \mathcal{H}_t)}{Q^j(\tau > T_j | \mathcal{H}_t)} \;-\; 1 \right).
\end{equation}
\textit{Says:} a one-period CDS spread is a ratio of forward
conditional survival probabilities under the $T_j$-forward measure,
with default-free discount factors completely cancelled --- analogous
to $F_k(t) = (P(t, T_{k-1})/P(t, T_k) - 1)/\alpha_k$ for a Libor
forward. This is the key conceptual result: \textit{the CDS spread
is fundamentally a probability ratio, not a rate}.

\paragraph{Multi-period spread as weighted one-period combination}
(paper's eq.~10):
\begin{equation}\label{eq:cdsmm-multi}
  R_{a,b}(t) \;=\; \sum_{j=a+1}^{b} \frac{\alpha_j \bar P(t, T_j)}{\sum_i \alpha_i \bar P(t, T_i)}\, R_j(t).
\end{equation}
\textit{Says:} the term-structure CDS spread is a defaultable-PV01-weighted
average of one-period spreads, exact analogue of the swap rate as a
PV01-weighted combination of Libor forwards.

\paragraph{The CDS-LMM under default--rate independence}
(\S5.1 of the paper):
\begin{equation}\label{eq:cdsmm-lmm}
  \dd R_k(t) \;=\; \bar\sigma_k(t)\, R_k(t)\, \dd \bar V^k_k(t),
  \qquad
  \dd F_k(t) \;=\; \sigma_k(t)\, F_k(t)\, \dd Z^k_k(t),
  \qquad
  \dd V_i \dd Z_j \;=\; 0,
\end{equation}
each rate a lognormal martingale under its own measure $\bar Q^k$
(for CDS) or $Q^k$ (for Libor), with vanishing cross-correlation.
The dynamics of $R_i$ under a foreign measure $\bar Q^j$ (with $i \ne j$):
\begin{equation}\label{eq:cdsmm-drift}
  \dd R_i(t) \;=\; \bar\sigma_i(t)\, R_i(t)
                \sum_{h=(i \wedge j)+1}^{i \vee j} h^j_i\, \rho_{i,h}\,
                \frac{\bar\sigma_h(t)\, R_h(t)}{R_h(t) + \alpha_h/\mathrm{LGD}}\, \dd t
                \;+\; \bar\sigma_i(t)\, R_i(t)\, \dd V^j_i(t),
\end{equation}
with $h^j_i = \Ind_{j<i} - \Ind_{j>i}$. \textit{Says:} the LMM
machinery transplants intact to CDS spreads, modulo a factor
$R_h + \alpha_h/\mathrm{LGD}$ in the drift that comes from the
defaultable bond ratio $\bar P(t, T_{h-1})/\bar P(t, T_h)$.

\paragraph{CMCDS pricing under drift-freezing} (eq.~at p.~21):
\begin{equation}\label{eq:cdsmm-cmcds}
  \text{CMCDS}_{a,b,c}(0)
  \;=\; \sum_{j=a+1}^{b} \alpha_j\, \bar P(0, T_j)
        \left[ \sum_{i=j}^{j+c} \bar w_{ij}(0)\, R_i(0)\, e^{T_{j-1}\, \sigma_i \sum_{h=j+1}^{i} \rho_{i,h}\, \sigma_h R_h(0)/(R_h(0)+\alpha_h/\mathrm{LGD})}
                \;-\; R_j(0) \right],
\end{equation}
with $\bar w_{ij}(t) = \alpha_i \bar P(t, T_i)/\sum_{h=j}^{j+c} \alpha_h \bar P(t, T_h)$.
\textit{Says:} closed-form CMCDS valuation under (i) the
drift-freezing approximation $\mu^j_i(R(t), t) \approx \mu^j_i(R(0), t)$,
(ii) representation of the multi-period swap rate as a linear
combination of one-period rates with PV01 weights, and (iii)
deterministic vols and constant correlations.

% =====================================================================
\section{Derivations \& commentary}

\begin{commentary}
The headline trick of the paper is the subfiltration $\mathcal{H}_t$
of default-free information, importing Jeanblanc-Rutkowski (2000)
into a Jamshidian (2004)-style market-model construction. The
problem in Schonbucher (1999) was that the defaultable-bond
numeraire $\Ind_{\tau > t}\bar B_k(t)$ is not strictly positive --- it
collapses to zero on the default event --- so the change-of-measure
machinery only gives a measure absolutely continuous with $Q$, not
equivalent. The paper's $C_{a,b}$ numeraire is the pre-default
defaultable PV01, which is strictly positive everywhere; under
$\mathcal{H}_t$-conditioning $1/C_{a,b}(t)$ is bounded; the resulting
$\bar Q^{a,b}$ is fully equivalent to $Q$.
\end{commentary}

\begin{commentary}
The conditional-independence assumption (Remark~4, paper p.~9) is
non-trivial: it says that \emph{the price of an $\mathcal{H}$-measurable
payoff depends only on $\mathcal{H}_t$, not on whether default has
already occurred}. Jamshidian (2004) notes this ``degrades the role
of the subfiltration''. The financial reading: pre-default valuations
of default-free claims don't depend on the fact of survival ---
sensible. With this, the Brigo (2005) Cox-process formula
generalises to the full subfiltration framework.
\end{commentary}

\begin{commentary}
\textbf{Why the term-structure model is hard.} Eq.~\eqref{eq:cdsmm-onerate}
expresses $R_j$ as a ratio of forward survival probabilities under
$Q^j$. So the defaultable bond \emph{ratio} $\bar P(t, T_{j-1})/\bar P(t, T_j)$
involves $Q^{j-1}(\tau > T_{j-1}|\mathcal{H}_t)/Q^j(\tau > T_j|\mathcal{H}_t)$,
a \emph{cross-measure} probability ratio --- not the same as
$Q^j(\tau > T_{j-1}|\mathcal{H}_t)/Q^j(\tau > T_j|\mathcal{H}_t)$ that
appears in the definition of $R_j$. Hence $\bar P(t,T_{j-1})/\bar P(t,T_j)$
cannot be expressed purely in terms of $R_j(t)$'s. This is the
``measure mismatch'' problem (paper p.~16).
\end{commentary}

\begin{commentary}
The mismatch disappears under independence of default and rates:
then the survival probabilities are the same under all forward
measures, and $\bar P(t, T_{j-1})/\bar P(t, T_j) = (1 + \alpha_j R_j(t)/\mathrm{LGD})(1 + \alpha_j F_j(t))$
(paper p.~18). This finally allows a tractable CDS-LMM. The price is
giving up the joint stochasticity of default-free interest rates and
credit spreads --- a major modelling concession, but acceptable in
practice for vanilla CDS-derivative pricing.
\end{commentary}

\begin{commentary}
\textbf{Empirical takeaways.} Table on p.~11 shows DT, DCX, FT CDS
options as of 26-Mar-2004: implied vols 50--55\%, an order of
magnitude higher than interest-rate swaptions but consistent with
Schonbucher (2004) historical estimates. Sensitivity to recovery
rate is small ($\pm 1$\%); sensitivity to a $\pm 50$~bp zero-curve
shift is also small ($\pm 0.5$--$0.7\%$). \emph{The implied vols are
the right way to compare CDS options across strikes and maturities,
not raw prices.}
\end{commentary}

\begin{commentary}
The CMCDS pricing (Sec.~5.2) is the killer application: a CMCDS pays
the running $(c+1)$-period CDS spread at each coupon date, which
combines a credit underlying with a maturity-rolling feature.
Without a CDS term-structure model, this is hard; with the LMM-style
construction here, it falls out via drift-freezing in closed form
\eqref{eq:cdsmm-cmcds}. The empirical example on Fiat shows the
``convexity adjustment'' (the difference between the naive valuation
ignoring CDS-rate dynamics and the full model) is meaningful at
typical vols and correlations.
\end{commentary}

% =====================================================================
\section{Validation}

\begin{validation}
This is a \emph{specification} of mathematical objects the paper
defines and the numerical values they should take. It says nothing
about how those objects are implemented or invoked.

\textbf{Quantities the paper defines:}
\begin{itemize}[noitemsep]
  \item Pre-default defaultable bond $\bar P(t, T)$ and the par CDS
        spread \eqref{eq:cdsmm-par}.
  \item Defaultable PV01 $C_{a,b}(t)$ \eqref{eq:cdsmm-pv01} and the
        corresponding measure $\bar Q^{a,b}$.
  \item Standard CDS option Black formula \eqref{eq:cdsmm-black}.
  \item One-period CDS spread $R_j(t)$ as a forward survival-probability
        ratio \eqref{eq:cdsmm-onerate}.
  \item Multi-period spread as weighted combination
        \eqref{eq:cdsmm-multi}.
  \item CDS-LMM dynamics \eqref{eq:cdsmm-lmm} and the drift under a
        foreign measure \eqref{eq:cdsmm-drift}.
  \item Closed-form CMCDS price \eqref{eq:cdsmm-cmcds} under
        drift-freezing.
\end{itemize}

\textbf{Canonical numerical cases:}
\begin{itemize}[noitemsep]
  \item European CDS options on 3 corporates (Deutsche Telecom $C_1$,
        Daimler Chrysler $C_2$, France Telecom $C_3$) on
        26-March-2004, recovery $\Rec = 0.4$, expiry $T_a$ = 20-Jun-2004
        (and $T_a' = $ 20-Dec-2004 for $C_1$), maturity $T_b$ =
        20-Jun-2009.
  \item Fiat CMCDS with $a = 0$, $b = 20$, $c = 20$ (resetting
        quarterly) on 20-Dec-2004, recovery $\Rec = 0.4$.
\end{itemize}

\textbf{Expected numerical outputs.}

\smallskip
\textit{CDS options} (paper p.~11, mid prices and implied vols):
\begin{tabular}{@{}lrrrr@{}}
\toprule
                & $R_{a,b}(0)$ (bp) & $K$ (bp) & Mid quote (bp) & Implied $\bar\sigma_{a,b}$ \\
\midrule
Option $C_1$ on $T_a$  & 61    & 60 & 32.5 & $62.16\%$ \\
Option $C_2$ on $T_a$  & 97.3  & 94 & 39   & $54.68\%$ \\
Option $C_3$ on $T_a$  & 62.7  & 61 & 25   & $52.01\%$ \\
Option $C_1$ on $T_a'$ & 65.4  & 61 & 35   & $51.45\%$ \\
\bottomrule
\end{tabular}
\smallskip

\textit{Recovery-rate sensitivity} ($C_1$ on $T_a$, table p.~12):
implied vol $50.02\%$ at $\Rec=20\%$, $50.31\%$ at $\Rec=40\%$,
$50.90\%$ at $\Rec=60\%$. \emph{Very weak dependence.}

\smallskip
\textit{Curve-shift sensitivity} ($\pm 50$~bp shift):
$C_1$ on $T_a$: vol $49.68\%$ at $-50$~bp, $50.31\%$ at $0$,
$50.93\%$ at $+50$~bp.

\smallskip
\textit{CMCDS Fiat convexity adjustment} (table p.~22, $a=0, b=20, c=20$):
\begin{tabular}{@{}lrrrr@{}}
\toprule
                                         & $\rho=0.7$ & $\rho=0.8$ & $\rho=0.9$ & $\rho=0.99$ \\
\midrule
$\sigma=0.1$                             & $0.000659$ & $0.000754$ & $0.000848$ & $0.000933$ \\
$\sigma=0.2$                             & $0.002662$ & $0.003047$ & $0.003435$ & $0.003784$ \\
$\sigma=0.4$                             & $0.011066$ & $0.012742$ & $0.014442$ & $0.015995$ \\
$\sigma=0.6$                             & $0.026619$ & $0.030964$ & $0.035464$ & $0.039652$ \\
\bottomrule
\end{tabular}
\smallskip

\textit{CMCDS participation rate} $\phi_{0,20,20}(\sigma, \rho)$
(decreasing in both vol and correlation): from $\sim 0.714$ at low
vol to $\sim 0.599$ at $\sigma = 0.6, \rho = 0.99$.

\textbf{Equation $\leftrightarrow$ quantity map:}

\smallskip
\begin{tabular}{@{}p{3.0cm}p{4.0cm}p{6.5cm}@{}}
\toprule
Paper eq.                                 & Symbol                          & Quantity \\
\midrule
\eqref{eq:cdsmm-payoff}                   & $\text{CDS}\Pi_t$               & CDS discounted payoff (tractable convention) \\
\eqref{eq:cdsmm-par}                      & $R_{a,b}(t)$                    & par CDS spread \\
\eqref{eq:cdsmm-pv01}                     & $C_{a,b}(t)$                    & defaultable PV01 (numeraire) \\
\eqref{eq:cdsmm-black}                    & $\text{CDSOption}_t$            & Black-formula CDS option price \\
\eqref{eq:cdsmm-onerate}                  & $R_j(t)$                        & one-period CDS as forward survival-prob ratio \\
\eqref{eq:cdsmm-multi}                    & $R_{a,b}$ in terms of $R_j$     & PV01-weighted aggregation \\
\eqref{eq:cdsmm-lmm}--\eqref{eq:cdsmm-drift} & LMM-style dynamics            & CDS spreads as lognormal martingales under natural measures \\
\eqref{eq:cdsmm-cmcds}                    & $\text{CMCDS}_{a,b,c}(0)$       & closed-form CMCDS under drift-freezing \\
\bottomrule
\end{tabular}
\smallskip

\textbf{Invariants and identities to verify:}
\begin{enumerate}[noitemsep]
  \item \emph{Defaultable bond limit}: as $\Rec \to 1$ (no loss),
        $R_{a,b}(t) \to 0$.
  \item \emph{Pre-default identity}: $\Ind_{\tau > t} \bar P(t, T)
        = \E^Q[D(t,T)\Ind_{\tau>T}|\mathcal{F}_t]$, so
        $\bar P(t, T) \le P(t, T)$ always.
  \item \emph{$C_{a,b}$ positivity}: under
        $Q(\tau > t|\mathcal{H}_t) > 0$ a.s.\ (paper's assumption),
        $C_{a,b}(t) > 0$ pre-default (Proposition~7 of paper).
  \item \emph{Telescoping of one-period spreads to multi-period}:
        \eqref{eq:cdsmm-multi} holds identically, by construction of
        $C_{a,b}$ from $\bar P(t, T_i)$.
  \item \emph{ATM CDS option}: at $K = R_{a,b}(t)$,
        $d_1 = -d_2 = \bar\sigma_{a,b}\sqrt{T_a-t}/2$, and the option
        price is approximately
        $R_{a,b}(t)\, C_{a,b}(t)\, \bar\sigma_{a,b}\sqrt{T_a-t}/\sqrt{2\pi}$.
  \item \emph{Implied vol back-out}: given a market quote
        (e.g.\ $C_1$: mid 32.5~bp on $R_{a,b}(0)=61$, $K=60$,
        $T_a - t \approx 0.24$Y), solving the Black formula
        \eqref{eq:cdsmm-black} for $\bar\sigma_{a,b}$ returns
        $62.16\%$.
  \item \emph{LMM consistency at $\rho = 0$}: in \eqref{eq:cdsmm-drift},
        all $\rho_{i,h} = 0$ collapses the drift to zero, recovering
        the martingale dynamics under $\bar Q^j = \bar Q^i$.
  \item \emph{CMCDS at $\sigma = 0$}: setting all
        $\bar\sigma_i \to 0$ in \eqref{eq:cdsmm-cmcds} collapses the
        exponential to 1 and gives the no-convexity-adjustment value
        $\sum \alpha_j \bar P(0, T_j) [\sum_i \bar w_{ij} R_i(0) - R_j(0)]$.
        Reproduces the $\rho = 0$ rows of the participation-rate
        table (since with $\sigma = 0$ the only correction is from
        the LMM-style replacement of multi-period rates).
  \item \emph{Conditional-independence test}: under independence of
        rates and default, the one-period spread becomes
        $R_j(t) = (\mathrm{LGD}/\alpha_j)\, [Q(\tau > T_{j-1}|\mathcal{H}_t)/Q(\tau > T_j|\mathcal{H}_t) - 1]$
        with $Q$ in place of $Q^j$. Verify this equals
        \eqref{eq:cdsmm-onerate} when $F_j$ and survival are uncorrelated.
\end{enumerate}
\end{validation}

% =====================================================================
\section{Glossary of symbols (paper's notation)}
\begin{tabular}{@{}p{3cm}p{12cm}@{}}
\toprule
$\tau$                      & Default time \\
$T_a, T_b$                  & Start and maturity of CDS premium leg \\
$\alpha_i$                  & Year fraction of period $(T_{i-1}, T_i]$ \\
$\mathrm{LGD}$              & Loss given default ($= 1 - $ recovery) \\
$R$                         & Running CDS premium (rate) paid by protection buyer \\
$R_{a,b}(t)$                & Par CDS spread for maturity $T_a \to T_b$ \\
$D(t, T)$                   & Stochastic bank-account discount $B_t/B_T$ \\
$P(t, T)$                   & Default-free zero-coupon bond price \\
$\bar P(t, T)$              & Pre-default defaultable bond price \\
$Q$                         & Risk-neutral measure (bank-account numeraire) \\
$Q^j$                       & $T_j$-forward measure under $P(t, T_j)$ \\
$\bar Q^{a,b}$              & Measure under defaultable-PV01 numeraire $C_{a,b}$ \\
$Q(\tau > t | \mathcal{H}_t)$ & Conditional survival probability (strictly positive) \\
$\mathcal{F}_t$             & Full filtration (default time observable) \\
$\mathcal{F}^\tau_t$        & Subfiltration generated by $\tau$ \\
$\mathcal{H}_t$             & Default-free subfiltration \\
$C_{a,b}(t)$                & Defaultable PV01 numeraire \\
$R_j(t) = R_{j-1,j}(t)$     & One-period forward CDS spread \\
$F_k(t)$                    & Default-free forward Libor (LMM variable) \\
$\bar\sigma_{a,b}$          & Implied volatility of $R_{a,b}$ under $\bar Q^{a,b}$ \\
$\rho_{i,h}$                & Correlation between Brownians of $R_i, R_h$ \\
$h^j_i$                     & Sign indicator: $\Ind_{j<i} - \Ind_{j>i}$ \\
$\bar w_{ij}(t)$            & PV01 weight $\alpha_i \bar P(t, T_i)/\sum_h \alpha_h \bar P(t, T_h)$ \\
$\text{Black}(F, K, v)$     & Standard Black formula \\
\bottomrule
\end{tabular}

% =====================================================================
\begin{todobox}
\begin{itemize}[noitemsep]
  \item Reproduce the implied-vol back-out on the four worked CDS
        options (Deutsche Telecom, Daimler Chrysler, France
        Telecom, with the $T_a, T_a'$ variants). Target: 62.16\%,
        54.68\%, 52.01\%, 51.45\% (paper p.~11).
  \item Verify the weak dependence of $\bar\sigma_{a,b}$ on recovery
        rate (variation $\le 0.5\%$ over $\Rec \in [0.2, 0.6]$) and
        the comparable weak dependence on parallel curve shifts.
  \item Implement the CDS-LMM under default--rate independence
        \eqref{eq:cdsmm-lmm} with vanishing cross-correlation;
        propagate one-period rates $R_k(t)$ and Libor rates $F_k(t)$
        in their own natural measures.
  \item Reproduce the Fiat CMCDS convexity-adjustment table
        ($a=0, b=20, c=20$, $\Rec=0.4$) across the
        $(\sigma, \rho)$ grid; target the exact values in
        \eqref{eq:cdsmm-cmcds}.
  \item Reproduce the participation-rate table (decreasing in
        $\sigma, \rho$).
  \item Implement the change-of-measure drift adjustment
        \eqref{eq:cdsmm-drift} and Monte-Carlo it against the
        drift-freezing closed form for a CMCDS; quantify the
        approximation error.
  \item Cross-check against Zhou (2008) (\texttt{zhou\_2008\_cds}):
        Brigo-Morini give the spread as a forward-survival-probability
        ratio (eq.~13), Zhou gives the bond-implied CDS as a
        weighted sum of bond coupon, Libor and bond discount; at
        consistent inputs they should agree on the par-spread level.
\end{itemize}
\end{todobox}

\end{document}
